% homework/hw04/hw04.tex
% Year 7 - Homework Packet 4
% Covers Maths 1.5 (tests for divisibility) and Maths 1.6 (square roots and
% cube roots). No science this week: the packet is two sheets, and the End of
% Unit 1 test is on Thursday, so the whole of it is maths revision.
%
% The spine of the packet is Section B, and it is the reason the two sections
% are in the same packet at all. A square number is a PAIR -- 36 = 6 x 6 --
% so every factor inside it turns up twice; a cube is three copies. That makes
% the divisibility tests from 1.5 into a rule-out test for 1.6:
%   * 3 divides but 9 does not  -> it cannot be a square number
%   * 2 divides but 4 does not  -> it cannot be a square number
%   * 2 divides but 8 does not  -> it cannot be a cube number
% and then into a method: the test that says "9 divides this" is exactly how
% you split 1296 into 9 x 144 and read off 3 x 12 = 36 without ever guessing.
%
% B1 deliberately contains 108, which passes every rule and is still not a
% square. The rules can only say NO. B1(g) makes the student say so.
%
% B3 question 3 closes the loop with Homework 3: 5832 was a row in that
% packet's tick grid. It passes every square rule here and is a cube.
%
% A1 is not decoration -- the two lists it makes the student write are the
% tool that Section B runs on, which is why the orange box above it does not
% print them.
%
% House style is shared: see homework/hw-style.tex. Method: the new-homework
% skill. Source is pure ASCII on purpose.
%
% Build:  pwsh .claude/skills/new-homework/scripts/build-hw.ps1 hw04

\documentclass[12pt,a4paper]{article}

% -- Answer key switch -------------------------------------------------------
\newif\ifkey
\ifdefined\TEACHER\keytrue\else
  \keyfalse        % \keyfalse = student copy   |   \keytrue = teacher copy
\fi

% -- House style -------------------------------------------------------------
% homework/hw-style.tex
% Shared house style for every Year 7 homework packet. \input this from the
% packet file, right after \documentclass and the answer-key switch.
%
% Do not fork this file per packet. If a packet needs a different look, the
% question is whether every packet should have it.
%
% The choices here are deliberate and were arrived at by printing and looking:
%   * No solid dark banners. Every header is a light tint with a coloured spine
%     on the left, so colour reads clearly without flooding a school printer.
%   * No marks and no mark totals. These are practice, not exam papers.
%   * Lato at 12pt with open leading, plus mathastext so inline maths is Lato
%     too. Without mathastext every "-6 + 4" is Computer Modern serif and the
%     sheet looks half-typeset.
%   * Writing lines are spaced for Year 7 handwriting, not adult margin notes.
%
% Provides: \wlines \blank \tickbox \sep \qmk-free marks (none) \hwsection
% \qhead \expr, the workspace box, and the remember / wordhelp / activity /
% yourturn coloured boxes. Palette matches the lesson decks exactly.

% -- Packages ----------------------------------------------------------------
\usepackage[T1]{fontenc}
\usepackage[utf8]{inputenc}
\usepackage[default]{lato}
\usepackage[a4paper,top=1.7cm,bottom=1.7cm,left=2.0cm,right=2.0cm,headheight=15pt]{geometry}
\usepackage{amsmath,amssymb}
\usepackage{xcolor}
\usepackage[most]{tcolorbox}
\usepackage{enumitem}
\usepackage{tikz}
\usetikzlibrary{arrows.meta}
\usepackage{array}
\usepackage{tabularx}
\usepackage{fancyhdr}
\usepackage{multicol}
\usepackage{needspace}
\usepackage{graphicx}

% Make maths use Lato too. Without this, every "-6 + 4" on the page is set in
% Computer Modern serif and the sheet looks half-typeset. Must come last.
\usepackage[italic]{mathastext}

\linespread{1.06}\selectfont

% -- The Learner's Book palette (same hex values as the lesson decks) ---------
\definecolor{cbteal}{HTML}{0087A8}
\definecolor{cbpurple}{HTML}{5C2483}
\definecolor{cborange}{HTML}{C25E12}
\definecolor{cbgreen}{HTML}{4A8B23}
\definecolor{cbred}{HTML}{C8102E}
\definecolor{cbblue}{HTML}{1A5FA8}
\definecolor{rulegray}{HTML}{9AA5AD}

% -- Answer lines and blanks -------------------------------------------------
% \wlines{n} draws n ruled writing lines, spaced wide enough for Year 7
% handwriting rather than for an adult with a fine pen.
\makeatletter
\newcommand{\wlines}[1]{%
  \par\vspace{0.75em}%
  \@tempcnta=0\relax
  \@whilenum\@tempcnta<#1\do{%
    \hbox to \linewidth{\color{rulegray}\leaders\hrule height 0.5pt\hfill}%
    \vspace{1.5em}%
    \advance\@tempcnta by 1\relax}%
  \par\vspace{-0.8em}%
}
\makeatother

% \blank[width] is an inline gap to write a short answer in.
\newcommand{\blank}[1][2.6cm]{\underline{\hspace{#1}}}

% A boxed space for showing working.
\newtcolorbox{workspace}[1][3.4cm]{%
  colback=white, colframe=rulegray, boxrule=0.5pt, arc=5pt,
  height=#1, left=6pt, right=6pt, top=4pt, bottom=4pt,
  before skip=9pt, after skip=11pt}

% An empty tick box.
\newcommand{\tickbox}{\fbox{\rule{0pt}{1.15em}\rule{1.15em}{0pt}}}

% A middot separator.
\newcommand{\sep}{\,\textperiodcentered\,}

% -- Section and question headings -------------------------------------------
% Light tint plus a fat coloured spine on the left. No solid dark banner: the
% colour reads clearly but barely costs the printer anything.
% needspace stops a heading stranding alone at the foot of a page. Kept modest:
% too large and it pushes headings over, leaving a worse hole than it prevents.
\newcommand{\hwsection}[2][cbteal]{%
  \needspace{8\baselineskip}%
  \par\vspace{13pt}%
  \noindent
  \begin{tcolorbox}[enhanced, nobeforeafter, width=\textwidth,
      colback=#1!7, colframe=#1!7, arc=5pt,
      borderline west={5pt}{0pt}{#1},
      left=14pt, right=10pt, top=7pt, bottom=7pt]
    {\large\bfseries\color{#1}#2}
  \end{tcolorbox}%
  \par\vspace{9pt}}

% \qhead[colour]{A2}{Moving along the line} -- a question heading.
\newcommand{\qhead}[3][cbteal]{%
  \needspace{6\baselineskip}%
  \par\vspace{11pt}%
  \noindent{\large\bfseries\color{#1}#2}\quad{\large\bfseries #3}\par\vspace{4pt}}

% -- Coloured boxes ----------------------------------------------------------
% Shared look: rounded, light fill, coloured rule, coloured title text on a
% pale title strip rather than a dark one.
\tcbset{hwbox/.style={
  enhanced, breakable, arc=6pt, boxrule=1pt,
  left=11pt, right=11pt, top=8pt, bottom=8pt,
  fonttitle=\bfseries, attach boxed title to top left={xshift=12pt, yshift=-8pt},
  boxed title style={arc=4pt, boxrule=1pt, left=7pt, right=7pt, top=2pt, bottom=2pt},
  top=13pt}}

% Orange: things to remember. Matches the deck's "Write This Down".
\newtcolorbox{remember}[1][Remember from the lesson]{hwbox,
  colback=cborange!6, colframe=cborange, coltitle=cborange,
  colbacktitle=white, title={#1}}

% Blue: English help. The barrier in this class is the wording, not the maths.
\newtcolorbox{wordhelp}[1][Word help]{hwbox,
  colback=cbblue!6, colframe=cbblue, coltitle=cbblue,
  colbacktitle=white, title={#1}}

% Purple: an activity or instruction.
\newtcolorbox{activity}[1][How to do this homework]{hwbox,
  colback=cbpurple!6, colframe=cbpurple, coltitle=cbpurple,
  colbacktitle=white, title={#1}}

% Green: your turn / a friendly nudge.
\newtcolorbox{yourturn}[1][Your turn]{hwbox,
  colback=cbgreen!7, colframe=cbgreen, coltitle=cbgreen!75!black,
  colbacktitle=white, title={#1}}

% -- Shorthands --------------------------------------------------------------
\newcommand{\dC}{\,^{\circ}\mathrm{C}}          % use inside maths: $-8\dC$
\newcommand{\key}[1]{{\color{cborange}\textbf{#1}}}
\newcommand{\eng}[1]{{\color{cbblue}\textbf{#1}}}

\setlist[enumerate]{itemsep=8pt, topsep=5pt, leftmargin=*}
\setlist[itemize]{itemsep=4pt, topsep=4pt, leftmargin=*}
\renewcommand{\arraystretch}{1.45}

% -- An expression the student is asked to work from -------------------------
\newcommand{\expr}[1]{%
  \tcbox[on line, colback=cbgreen!10, colframe=cbgreen!55, boxrule=1pt, arc=4pt,
         left=9pt, right=9pt, top=4pt, bottom=4pt]{\large\bfseries $#1$}}

% -- Running head and foot ---------------------------------------------------
% The packet overrides \fancyhead[L] with its own title.
\pagestyle{fancy}
\fancyhf{}
\fancyhead[L]{\small\color{black!55}Year 7\sep Homework}
\fancyhead[R]{\small\color{black!55}Name: \blank[3.4cm]}
\fancyfoot[C]{\small\color{black!55}Page \thepage}
\renewcommand{\headrulewidth}{0pt}
\renewcommand{\headrule}{\vspace{2pt}\hbox to\headwidth{\color{cbteal!45}\leaders\hrule height 1.2pt\hfill}}



% -- Per-packet running head -------------------------------------------------
\fancyhead[L]{\small\color{black!55}Year 7\sep Homework 4}

\begin{document}

% ===========================================================================
% COVER BLOCK
% ===========================================================================
\thispagestyle{fancy}

\begin{tcolorbox}[enhanced, colback=cbpurple!7, colframe=cbpurple!7, arc=7pt,
                  borderline west={6pt}{0pt}{cbpurple},
                  left=18pt, right=14pt, top=12pt, bottom=12pt]
  {\color{cbpurple}\bfseries YEAR 7\sep UNIT 1\sep HOMEWORK 4}\\[6pt]
  {\color{cbpurple}\Huge\bfseries Squares, Cubes and Roots}\\[10pt]
  {\color{black!70}Mathematics 1.6 --- Square Roots and Cube Roots\sep
   1.5 --- Tests for Divisibility}
\end{tcolorbox}

\vspace{8pt}
\noindent
\begin{tabularx}{\textwidth}{@{}X X@{}}
  \textbf{Name:} \blank[5.6cm] & \textbf{Class:} \blank[5.6cm] \\
  \textbf{Date given:} \blank[4.6cm] & \textbf{Date due:} \blank[5.0cm] \\
\end{tabularx}

\vspace{6pt}
\begin{activity}
  \begin{itemize}
    \item Write your answers \textbf{on this sheet}. \key{No calculator.}
    \item Always \key{write the calculation, not just the answer}.
    \item Do \textbf{A1} first. You need those two lists for the whole packet.
    \item In \textbf{Section B} you use last week's divisibility tests to find
          roots.
  \end{itemize}
\end{activity}

% ===========================================================================
% SECTION A -- MATHS 1.6
% ===========================================================================
\hwsection{Section A \textperiodcentered\ Mathematics 1.6 --- Squares and Cubes}

% This box carries the notation and the there-and-back idea only. It must NOT
% print the two lists: A1 is the student writing them out.
\begin{remember}
  \key{$5^2$} means $5 \times 5 = 25$. We say \key{``five squared''}. The small 2
  counts \textbf{how many fives}, so $5^2$ is \key{not} $5 \times 2$.
  \key{$5^3$} means $5 \times 5 \times 5 = 125$. We say \key{``five cubed''}.

  \smallskip
  A \key{root} goes \key{backwards}. $5^2 = 25$, so the \key{square root} of 25 is
  5. $5^3 = 125$, so the \key{cube root} of 125 is 5.
\end{remember}

% -- A1 ---------------------------------------------------------------------
\qhead{A1}{The two lists}

\noindent\textbf{(a)} Fill in the \textbf{twelve square numbers}. Use your notebook.

\vspace{5pt}
\noindent
\begin{tabularx}{\textwidth}{|c|*{12}{>{\centering\arraybackslash}X|}}
  \hline
  $n$ & 1 & 2 & 3 & 4 & 5 & 6 & 7 & 8 & 9 & 10 & 11 & 12 \\
  \hline
  $n^2$ \rule[-0.8em]{0pt}{2.3em} & & & & & & & & & & & & \\
  \hline
\end{tabularx}

\vspace{10pt}
\noindent\textbf{(b)} Now the \textbf{cube numbers}.

\vspace{5pt}
\noindent
\begin{tabularx}{\textwidth}{|c|*{7}{>{\centering\arraybackslash}X|}}
  \hline
  $n$ & 1 & 2 & 3 & 4 & 5 & 6 & 10 \\
  \hline
  $n^3$ \rule[-0.8em]{0pt}{2.3em} & & & & & & & \\
  \hline
\end{tabularx}

\vspace{9pt}
\noindent\textbf{(c)} One number is on \textbf{both} lists. Which one?
\blank[3.2cm]

% -- A2 ---------------------------------------------------------------------
\qhead{A2}{Square it, cube it, root it}

\noindent\textbf{(a)} Work out \textbf{both}. Then write which one is bigger.

\vspace{4pt}
\noindent\hspace*{1em}$7^2 = \blank[2.2cm]$ \hspace{1.4cm}
$7 \times 2 = \blank[2.2cm]$ \hspace{1.4cm} Bigger: \blank[2.4cm]

\vspace{10pt}
\noindent\textbf{(b)} Use your lists from A1.

\vspace{6pt}
\begin{multicols}{3}
\begin{enumerate}[label=(\roman*), itemsep=9pt]
  \item $11^2 = \blank[1.5cm]$
  \item $12^2 = \blank[1.5cm]$
  \item $4^3 = \blank[1.5cm]$
  \item $10^3 = \blank[1.5cm]$
  \item $\sqrt{81} = \blank[1.5cm]$
  \item $\sqrt{144} = \blank[1.5cm]$
  \item $\sqrt[3]{27} = \blank[1.5cm]$
  \item $\sqrt[3]{216} = \blank[1.5cm]$
  \item $\sqrt[3]{64} - \sqrt{100} = \blank[1.5cm]$
\end{enumerate}
\end{multicols}

\vspace{2pt}
\noindent\textbf{(c)} One answer in (b) is \textbf{negative}. Which one?
\blank[3.2cm]

% -- A3 ---------------------------------------------------------------------
\qhead{A3}{Say it, then write it}

\begin{wordhelp}[Word help --- these words tell you what to do]
  {\renewcommand{\arraystretch}{1.16}%
  \begin{tabularx}{\linewidth}{@{}l@{\hspace{1.4em}}X@{}}
    \eng{square 9}              & $9^2 = 9 \times 9$ \\
    \eng{the square of 9}       & the same thing: $9^2 = 81$ \\
    \eng{the square root of 81} & 9, because $9 \times 9 = 81$ \\
    \eng{cube 4}                & $4^3 = 4 \times 4 \times 4$ \\
    \eng{the cube root of 64}   & 4, because $4 \times 4 \times 4 = 64$ \\
    \eng{a square number}       & a number on your first list \\
    \eng{consecutive}           & one after another, nothing missed out. 6 and 7 are
                                  consecutive. 36 and 49 are consecutive
                                  \textbf{square} numbers. \\
  \end{tabularx}}
\end{wordhelp}

\noindent Write the calculation for each sentence, then work it out.
\key{Calculation first.}

% The calculation and answer go on their own line under the sentence. With
% \hfill on the same line, the long sentences in 2 and 3 wrapped and left the
% Answer blank stranded on a line of its own.
\vspace{2pt}
\begin{enumerate}[label=\arabic*., itemsep=5pt]
  \item Square 11. \\[3pt]
        Calculation: \blank[6.0cm] \quad Answer: \blank[2.4cm]
  \item Find the square root of 121. \\[3pt]
        Calculation: \blank[6.0cm] \quad Answer: \blank[2.4cm]
  \item Subtract the cube root of 27 from the square of 8. \\[3pt]
        Calculation: \blank[6.0cm] \quad Answer: \blank[2.4cm]
\end{enumerate}

% -- A4 ---------------------------------------------------------------------
\qhead[cbgreen]{A4}{Trap the root}

% Kept to two lines on purpose. At three it broke over the page and left the
% words "whole numbers." stranded alone at the top of the next one.
\begin{yourturn}[When the root is not a whole number]
  45 is not a square number, so $\sqrt{45}$ is not whole. But 45 is between
  \key{36} and \key{49}, and those two \textbf{are} on the list. So $\sqrt{45}$ is
  between \key{6} and \key{7} --- two \key{consecutive} whole numbers.
\end{yourturn}

\noindent Write the two consecutive whole numbers each root is between.

\vspace{8pt}
\begin{multicols}{2}
\begin{enumerate}[label=(\alph*), itemsep=11pt]
  \item $\blank[1.1cm] < \sqrt{30} < \blank[1.1cm]$
  \item $\blank[1.1cm] < \sqrt{70} < \blank[1.1cm]$
  \item $\blank[1.1cm] < \sqrt{110} < \blank[1.1cm]$
  \item $\blank[1.1cm] < \sqrt[3]{50} < \blank[1.1cm]$
\end{enumerate}
\end{multicols}

\vspace{8pt}
\noindent\textbf{(e)} Two \textbf{consecutive square numbers} add up to \textbf{85}.
Which two? \blank[4.6cm]

\vspace{6pt}
\noindent\hspace*{1em}Their square roots are \blank[1.6cm] and \blank[1.6cm]

% ===========================================================================
% SECTION B -- 1.5 MEETS 1.6
% ===========================================================================
\hwsection{Section B \textperiodcentered\ Maths 1.5 and 1.6 --- Using the Tests on a Root}

% "Screen" was the word here in the first draft. It is hard English for this
% class and it is not in the book, so the rules are just called rules.
\begin{remember}[A square is a pair. A cube is three copies.]
  $36 = 6 \times 6$. Both numbers in the pair are the \textbf{same}, so every
  factor inside one is inside the other as well. If \key{3} divides a square
  number, then \key{9} divides it too.

  \smallskip
  A cube is \key{three} copies: $8 = 2 \times 2 \times 2$. So a factor is in there
  \textbf{three times}. If \key{2} divides a cube number, then \key{8} divides it
  too.

  \smallskip
  {\renewcommand{\arraystretch}{1.14}%
  \begin{tabularx}{\linewidth}{@{}l@{\hspace{1.4em}}X@{}}
    \key{Rule 1} & divides by 3 but \textbf{not} by 9 $\Rightarrow$ \textbf{not} a
                   square number \\
    \key{Rule 2} & divides by 2 but \textbf{not} by 4 $\Rightarrow$ \textbf{not} a
                   square number \\
    \key{Rule 3} & divides by 2 but \textbf{not} by 8 $\Rightarrow$ \textbf{not} a
                   cube number \\
  \end{tabularx}}

  \smallskip
  \key{Careful.} These rules can only say \key{NO}. If a number passes, that does
  \textbf{not} prove it is a square.
\end{remember}

% -- B1 ---------------------------------------------------------------------
\qhead{B1}{Can it be a square number?}

\noindent Use the tests and fill in every box. In the last column write \textbf{NO}
if a rule says no, or \textbf{maybe} if it passes. If the number is \textbf{odd},
put a dash in the \textbf{by 2} and \textbf{by 4} columns.

\vspace{8pt}
\noindent
\begin{tabularx}{\textwidth}{|p{1.7cm}|p{1.9cm}|>{\centering\arraybackslash}X|>{\centering\arraybackslash}X|>{\centering\arraybackslash}X|>{\centering\arraybackslash}X|p{2.6cm}|}
  \hline
  \textbf{Number} & \textbf{Digit sum} & \textbf{by 3?} & \textbf{by 9?} &
  \textbf{by 2?} & \textbf{by 4?} & \textbf{Square?} \\
  \hline
  % 300 used to be a sixth row. It is ruled out by the same rule as 132
  % (divides by 3, not by 9), so it cost a centimetre and taught nothing new.
  \textbf{132}  \rule[-0.85em]{0pt}{2.3em} & & & & & & \\ \hline
  \textbf{441}  \rule[-0.85em]{0pt}{2.3em} & & & & & & \\ \hline
  \textbf{162}  \rule[-0.85em]{0pt}{2.3em} & & & & & & \\ \hline
  \textbf{1296} \rule[-0.85em]{0pt}{2.3em} & & & & & & \\ \hline
  \textbf{108}  \rule[-0.85em]{0pt}{2.3em} & & & & & & \\ \hline
\end{tabularx}

\vspace{9pt}
% Was an abstract "write a sentence about what these rules can and cannot
% prove" question -- too hard, and not concrete maths. Replaced with the
% trap-the-root check from A4, a skill they already have.
\noindent\textbf{(g)} One number passed \textbf{every} rule. Which one?
\blank[3.0cm]

\vspace{6pt}
\noindent\hspace*{1em}Check it against your A1 list. It is between
\blank[1.4cm] and \blank[1.4cm] --- two square numbers.

\vspace{6pt}
\noindent\hspace*{1em}So is it a square number? \blank[2.4cm] (yes or no)

% -- B2 ---------------------------------------------------------------------
\qhead{B2}{Split it, then find the root}

\begin{remember}[Split the number into two pieces you know]
  Split the number into two pieces. Find the root of each piece. Then
  \key{multiply the two roots}. \key{Good pieces:} 9, 4 and 100 are squares (roots
  3, 2, 10); 8 and 1000 are cubes (roots 2 and 10).

  \smallskip
  \key{Example.} $\sqrt{576}$. The digits add to $5 + 7 + 6 = 18$, so \key{9
  divides it}. $576 \div 9 = 64$. Now $\sqrt{9} = 3$ and $\sqrt{64} = 8$, so
  $\sqrt{576} = 3 \times 8 = \mathbf{24}$. Check: $24 \times 24 = 576$.
\end{remember}

% The "Divide by" column used to be blank -- find the right test yourself, THEN
% do the split. That was two hard steps in a row. Now the divisor is given, so
% the only jobs left are the division, the roots, and the multiplication.
\noindent Now do these. The number to divide by is given. Fill in the rest.

\vspace{8pt}
\noindent
\begin{tabularx}{\textwidth}{|p{2.2cm}|p{1.8cm}|X|X|p{1.9cm}|}
  \hline
  \textbf{Find} & \textbf{Divide by} & \textbf{Split it} &
  \textbf{The two roots} & \textbf{Answer} \\
  \hline
  % 441 and 1296 are the two rows B1 left as "maybe" -- this is where they get
  % settled, so neither of them can be the row that goes if space runs short.
  $\sqrt{441}$    \rule[-0.75em]{0pt}{2.1em} & 9    & & & \\ \hline
  $\sqrt{1296}$   \rule[-0.75em]{0pt}{2.1em} & 9    & & & \\ \hline
  $\sqrt[3]{1728}$\rule[-0.75em]{0pt}{2.1em} & 8    & & & \\ \hline
  $\sqrt[3]{8000}$\rule[-0.75em]{0pt}{2.1em} & 1000 & & & \\ \hline
\end{tabularx}

% -- B3 ---------------------------------------------------------------------
% B3 and B4 are in this order on purpose, and it is what took the packet from
% five pages to four. B3 is the big block that cannot break -- a box, four
% steps and a workspace -- so wherever it sat mid-section it stranded a third
% of a page. Last in the section it opened a hole; here it does not, and the
% leftover space at the foot of the packet lands on B4's workspace instead of
% on nothing. It also puts the deadpan line last, which is where it belongs.
\qhead[cbgreen]{B3}{The challenge: 5832}

\begin{yourturn}[5832, from Homework 3]
  In Homework 3 you ticked \key{5832} for 2, 3, 4, 6, 8 and 9. So it passes
  \key{every} rule in B1 --- but it is \key{not} a square number. It \key{is} a
  cube number. Find its cube root.
\end{yourturn}

\noindent\textbf{Step 1.} 5832 divides by 8. Do the division.
\eng{Tip: dividing by 8 is the same as halving three times.}
\quad $5832 \div 8 = \blank[3.0cm]$

\vspace{6pt}
\noindent\textbf{Step 2.} Your answer to Step 1 is $9 \times 9 \times 9$, so its
cube root is \blank[2.2cm]

\vspace{6pt}
\noindent\textbf{Step 3.} $\sqrt[3]{8} = \blank[1.6cm]$, so
$\sqrt[3]{5832} = \blank[1.6cm] \times \blank[1.6cm] = \blank[2.0cm]$

\vspace{6pt}
% Was one open box asking students to "check by multiplying by itself twice" --
% they had to invent their own two-step structure. Now the two multiplications
% are given as separate blanks.
\noindent\textbf{Step 4.} Check your answer.
\quad $18 \times 18 = \blank[2.2cm]$ \quad then \quad
$\blank[1.6cm] \times 18 = \blank[2.4cm]$

% -- B4 ---------------------------------------------------------------------
% One word problem, not two, and it is a SQUARE on purpose: cubes already have
% two rows of B2 and the whole of B3, and this is the only place in the packet
% where a square root has to be dug out of an English sentence. 400 also
% exercises the one useful split that B2 never uses -- the test for 4.
\qhead{B4}{Word problem}

\noindent\textbf{The photograph.} Mr Bowen puts \textbf{400 students} in a
\textbf{perfect square}. The photographer climbs onto the roof. How many students
are in each row? Which test did you use to split 400?

\begin{workspace}[2.8cm]\end{workspace}

% ===========================================================================
% ANSWER KEY (teacher copy only)
% ===========================================================================
\ifkey
\newpage
\hwsection[cbred]{Answer Key --- Teacher Copy}

% The key is for an adult at a desk, not a 12-year-old writing on it.
\linespread{1.0}\selectfont
\setlist[enumerate]{itemsep=2pt, topsep=3pt, leftmargin=*}
\setlist[itemize]{itemsep=2pt, topsep=3pt, leftmargin=*}

\qhead[cbred]{A1}{The two lists}
\noindent (a) 1, 4, 9, 16, 25, 36, 49, 64, 81, 100, 121, 144. \\
(b) 1, 8, 27, 64, 125, 216, and $10^3 = 1000$. \\
(c) \textbf{64} --- $8 \times 8$ and $4 \times 4 \times 4$.

\smallskip
\noindent{\small Mark this one first and mark it hard. A student who cannot fill
these two rows cannot do Section B at all, because Section B is nothing but
splitting a big number into two entries from these lists. If more than a few of the
class leave gaps, re-copy the lists at the board before going through anything
else in the packet.}

\qhead[cbred]{A2}{Out, and back again}
\noindent (a) $7^2 = 49$, \ $7 \times 2 = 14$. \ $7^2$ is bigger. \\
(b) (i) 121 \quad (ii) 144 \quad (iii) 64 \quad (iv) 1000 \quad (v) 9 \quad
(vi) 12 \quad (vii) 3 \quad (viii) 6 \quad (ix) $4 - 10 = -6$ \\
(c) \textbf{(ix)}, which comes to $-6$.

\smallskip
\noindent{\small (a) is the whole reason this question exists. An answer of 14 for
$7^2$ is the small-2-means-times-2 error and it is worth ten seconds at the board,
not a red cross. In (b), (ix) is the one they will call wrong: the class has
decided a root question cannot produce a negative, and $4 - 10$ says otherwise ---
that is 1.1 arriving inside a 1.6 question. Part (c) exists to make them notice it
rather than quietly change it to 6.}

\qhead[cbred]{A3}{Say it, then write it}
\begin{enumerate}[label=\arabic*., itemsep=2pt]
  \item $11^2 = 11 \times 11 = 121$
  \item $\sqrt{121} = 11$
  \item $8^2 - \sqrt[3]{27} = 64 - 3 = 61$
\end{enumerate}
\noindent{\small Question 3 is the English question, not the arithmetic one.
\emph{Subtract A from B} means $B - A$, so the calculation is $64 - 3$. A student
who writes $3 - 64$ has read the sentence left to right and has made the same error
they made with subtraction in 1.1 --- send them back to re-read the sentence rather
than marking the arithmetic. Questions 1 and 2 are on the page next to each other
on purpose: same numbers, opposite directions.}

\qhead[cbred]{A4}{Trap the root}
\noindent (a) $5 < \sqrt{30} < 6$ \quad (b) $8 < \sqrt{70} < 9$ \quad
(c) $10 < \sqrt{110} < 11$ \quad (d) $3 < \sqrt[3]{50} < 4$ \\
(e) \textbf{36 and 49}. Their roots are \textbf{6} and \textbf{7}.

\smallskip
\noindent{\small (d) is the one to check, because it is the only cube: they need
$27 < 50 < 64$, and a student running on the square list will answer 7 and 8.
(e) is the \emph{consecutive} question. Consecutive square numbers are 11 apart
here, not 1 apart, and a student who reads consecutive as ``one apart'' will hunt
for two numbers differing by 1 and give up. If several fail it, that is a
vocabulary lesson for Monday, not a marking problem.}

\qhead[cbred]{B1}{Can it be a square number?}
\noindent
\begin{tabularx}{\linewidth}{@{}p{1.6cm}p{1.7cm}p{1.2cm}p{1.2cm}p{1.2cm}p{1.2cm}X@{}}
  \textbf{132}  & sum 6  & yes & no  & yes & no  & \textbf{NO} \\
  \textbf{441}  & sum 9  & yes & yes & --- & --- & maybe (it is $21^2$) \\
  \textbf{162}  & sum 9  & yes & yes & yes & no  & \textbf{NO} \\
  \textbf{1296} & sum 18 & yes & yes & yes & yes & maybe (it is $36^2$) \\
  \textbf{108}  & sum 9  & yes & yes & yes & yes & maybe --- but it is not \\
\end{tabularx}

\smallskip
\noindent (g) \textbf{108}, between \textbf{100} and \textbf{121} ($10^2$ and
$11^2$). Not a square number.

\smallskip
\noindent{\small The two rows that teach are \textbf{162} and \textbf{108}, and they
teach opposite things. 162 passes the 3-and-9 rule and is killed by the 2-and-4
one, so a student who only ran the digit sum gets it wrong --- both rules, every
time. 108 is the honest one: it passes everything and is still not a square, which
is the difference between a test and a proof. This used to be an open question
asking students to explain that in their own words, which was too abstract for the
class; it is now the same trap-the-root check as A4, applied to 108 instead of a
root. If a student answers ``yes'' here, send them back to count: $10 \times 10 =
100$ and $11 \times 11 = 121$, and 108 is neither. A ``maybe'' written for 441 and
1296 in the table is the correct answer even though both really are squares --- do
not accept ``yes'' for those two on the strength of the rules alone.}

\qhead[cbred]{B2}{Split it, then find the root}
\noindent
\begin{tabularx}{\linewidth}{@{}p{2.3cm}p{1.4cm}p{2.6cm}p{2.6cm}X@{}}
  $\sqrt{441}$     & $\div 9$    & $9 \times 49$   & 3 and 7  & \textbf{21} \\
  $\sqrt{1296}$    & $\div 9$    & $9 \times 144$  & 3 and 12 & \textbf{36} \\
  $\sqrt[3]{1728}$ & $\div 8$    & $8 \times 216$  & 2 and 6  & \textbf{12} \\
  $\sqrt[3]{8000}$ & $\div 1000$ & $1000 \times 8$ & 10 and 2 & \textbf{20} \\
\end{tabularx}

\smallskip
\noindent Divisions worth having to hand: $441 \div 9 = 49$, $1296 \div 9 = 144$,
$1728 \div 8 = 216$, $8000 \div 1000 = 8$.

\smallskip
\noindent{\small The divisor is now printed for every row, which used to be the
hardest part --- deciding whether to use 9, 8 or 1000 was one guess before the
maths even started, and that made the question feel like four separate puzzles
rather than one method practised four times. With the divisor given, this is the
question the packet was built around: $\sqrt{1296}$ shows whether it has landed,
since 1296 is nowhere near their A1 list and dividing by 9 is the only route to 36.
The error to expect is \textbf{adding} the roots instead of multiplying them ---
$3 + 12 = 15$; tell them to check by squaring, because $15^2 = 225$ and not 1296.}

\qhead[cbred]{B3}{The challenge: 5832}
\noindent Step 1: $5832 \div 8 = \mathbf{729}$ (halve three times: $5832 \to 2916
\to 1458 \to 729$). \\
Step 2: $\sqrt[3]{729} = \mathbf{9}$. \\
Step 3: $\sqrt[3]{8} = 2$, so $\sqrt[3]{5832} = 2 \times 9 = \mathbf{18}$. \\
Step 4: $18 \times 18 = \mathbf{324}$, and $324 \times 18 = \mathbf{5832}$.

\smallskip
\noindent{\small Worth doing at the board even if most of the class got it, because
it is the entire unit in one question: a divisibility test from 1.5 picks the split,
a division carries it out, a cube root from 1.6 finishes it, and the check in Step 4
is two two-digit multiplications --- which is exactly what the extra sheet is for.
The halving tip in Step 1 matters: $5832 \div 8$ by long division is the hardest
single step in the packet, and most students have not thought to use $\div 2, \div
2, \div 2$ instead. Step 4 is now two separate blanks rather than one open box, so
a student who gets 324 but stalls on $324 \times 18$ has still shown the harder
half of the check.

\smallskip
This question is also the answer to B1(g) arriving in a form nobody can argue with.
5832 passed every rule and is not a square, so if the class was unconvinced
earlier, come back to B1(g) after this rather than re-teaching it cold.}

\qhead[cbred]{B4}{Word problem}
\noindent $\sqrt{400} = 20$, so \textbf{20 students in each row}. The split: 400
ends in 00, so $400 = 4 \times 100$, and $\sqrt{4} = 2$, $\sqrt{100} = 10$, giving
$2 \times 10 = 20$.

\smallskip
\noindent{\small A student who simply knows that $20 \times 20 = 400$ has the right
answer --- but the question asked which test they used, so ask for the split as
well. This is the only place in the packet where a square root has to be dug out of
an English sentence, so mark the reading as carefully as the arithmetic: the phrase
doing the work is \emph{a perfect square}, and a student who multiplies $400 \times
400$ has not found it. The photographer on the roof is the deadpan line. Read it
straight and move on.}

\fi
\end{document}
